\section{Global near-diagonal triage of the Mersenne excess}
\label{sec:mersenne-near-diagonal-triage}

The localized gate of
Corollary~\ref{cor:localized-exact-mersenne-gate} identifies the moving
divisor range on which an obstruction can survive.  We now combine that
localization with the exact-order Wieferich decomposition.  The result does
not estimate the remaining mass, but it proves that a failure cannot be
distributed among unspecified cyclotomic effects: one of three explicit
mechanisms must carry a positive linear logarithmic mass.

For an odd prime $q$, put
\[
 d_q=\operatorname{ord}_q(2),\qquad
 w_q=v_q(2^{d_q}-1),
\]
and, for $d\ge1$, let
\[
 \mathcal W_d=\{q:d_q=d,\ w_q\ge2\},\qquad
 D_d=\prod_{\substack{d_q=d\\w_q\ge3}}q^{w_q-2}.
\]
Fix $\delta>0$ and set
\[
 Y_d=\frac{\varphi(d)^2}{\log\log(3d)}.
\]
The logarithm of the canonical excess block has the exact four-layer
decomposition
\begin{equation}\label{eq:four-layer-mersenne-block}
 \log E_d=S_d+T_{d,\delta}+X_{d,\delta}+G_d,
\end{equation}
where
\[
\begin{aligned}
 S_d&=\sum_{\substack{q\in\mathcal W_d\\q\le Y_d}}\log q,
&T_{d,\delta}&=\sum_{\substack{q\in\mathcal W_d\\
      Y_d<q\le d^{2+\delta}}}\log q,\\
 X_{d,\delta}&=\sum_{\substack{q\in\mathcal W_d\\
      q>d^{2+\delta}}}\log q,
&G_d&=\log D_d.
\end{aligned}
\]
Indeed, the exponent of $q$ in $E_d$ is
$w_q-1=1+(w_q-2)$.  Moreover, LTE gives
\[
 v_q(2^{q-1}-1)=v_q(2^{d_q}-1)=w_q,
\]
so $\mathcal W_d$ consists precisely of base-two Wieferich primes of exact
order $d$.

Put
\[
 H(m)=\{\log\log(3m)\}^2,
 \qquad
 \mathcal N_m=\{d\mid m:d>m e^{-H(m)}\}.
\]
The cyclotomic cap $E_d\le3^{\varphi(d)}$ and the exact logarithmic-deficit
moment prove
\begin{equation}\label{eq:far-mersenne-mass-small}
 \sum_{\substack{d\mid m\\d\notin\mathcal N_m}}\log E_d=o(m).
\end{equation}
The uniform Brun--Titchmarsh argument for primes $q\equiv1\pmod d$ below
$Y_d$ gives $S_d=o(\varphi(d))$.  Since the least member of
$\mathcal N_m$ tends to infinity and
$\sum_{d\mid m}\varphi(d)=m$, this pointwise uniform estimate sums to
\begin{equation}\label{eq:near-small-wieferich-mass}
                     \sum_{d\in\mathcal N_m}S_d=o(m).
\end{equation}

\begin{theorem}[Global three-arm obstruction]
\label{thm:global-mersenne-three-arm-obstruction}
Suppose that, for some $\epsilon>0$ and an unbounded sequence $m_j$,
\[
                 \sum_{d\mid m_j}\log E_d\ge\epsilon m_j.
\]
After passage to a subsequence, one fixed alternative holds for every
sufficiently large $j$:
\begin{enumerate}
\item
$\sum_{d\in\mathcal N_{m_j}}G_d\ge\epsilon m_j/6$;
\item there are at least
\[
 \frac{\epsilon m_j}{6(2+\delta)\log m_j}
\]
distinct base-two Wieferich primes $q$ such that
\[
 d_q\mid m_j,\quad d_q>m_je^{-H(m_j)},\quad
 Y_{d_q}<q\le d_q^{2+\delta};
\]
\item
$\sum_{d\in\mathcal N_{m_j}}X_{d,\delta}\ge\epsilon m_j/6$.
Every prime counted in the last alternative satisfies
$\operatorname{ord}_q(2)<q^{1/(2+\delta)}$.
\end{enumerate}
\end{theorem}

\begin{proof}
By \eqref{eq:far-mersenne-mass-small} and
\eqref{eq:near-small-wieferich-mass}, the far mass and the near-small mass
are each at most $\epsilon m_j/4$ after deleting finitely many indices.
Summing \eqref{eq:four-layer-mersenne-block} therefore gives
\[
 \sum_{d\in\mathcal N_{m_j}}
   (G_d+T_{d,\delta}+X_{d,\delta})\ge\frac{\epsilon m_j}{2}.
\]
One of the three sums is at least $\epsilon m_j/6$; infinite pigeonhole
fixes the same arm on a subsequence.  In the transition case each summand is
at most $(2+\delta)\log m_j$.  Exact-order prime sets for different $d$ are
disjoint, so division by this cap proves the asserted cardinality.  The
extreme inequality $q>d_q^{2+\delta}$ gives the final order bound.
\end{proof}

\begin{corollary}[Three sufficient aggregate gates]
\label{cor:three-aggregate-mersenne-gates}
For one fixed $\delta>0$, assume
\[
 \sum_{d\in\mathcal N_m}G_d=o(m),\qquad
 \#\{\text{transition primes at }m\}=o(m/\log m),
 \qquad
 \sum_{d\in\mathcal N_m}X_{d,\delta}=o(m).
\]
Then $\log W_m=o(m)$.
\end{corollary}

\begin{proof}
The transition sum is at most $(2+\delta)\log m$ times its support size.
Adding the three hypotheses to
\eqref{eq:far-mersenne-mass-small} and
\eqref{eq:near-small-wieferich-mass} yields
$\sum_{d\mid m}\log E_d=o(m)$.  The exact factorization
$W_m=L_m\prod_{d\mid m}E_d$, with $L_m\mid m$, finishes the proof.
\end{proof}

None of the three disjuncts can be deleted by positivity alone.  On a
one-point near set, take target mass one, far and small masses zero, and
respectively
\[
       (G,T,X)=(1,0,0),\quad(0,1,0),\quad(0,0,1).
\]
These are complete counterexamples to each corresponding two-disjunct
ordered-ring strengthening.  They are abstract mass ledgers rather than
Mersenne blocks, so they do not eliminate an arithmetic arm.

The companion Lean module checks the finite near/far subtraction, the
$1/6$ trichotomy, the transition-cardinality lower bound, the five-arm
closure inequality, and all three full-premise counterexamples.  Its
analytic inputs remain the proved and cited statements above; no
distribution assertion is inserted as a formal axiom.  Present results on
primitive cyclotomic factors and generalized Wieferich congruences
\cite{GranvillePrimitive2012,FelliniMurty2026,FalkHarringtonJones2026}
do not bound any one of the three weighted moving arms.  All three routes
therefore remain active.
